% GENERATED FILE. Edit canonical lessons, then run npm run generate:books.
% canonical-source-hash: fnv1a64:e7f518330aae51a6
% canonical-lessons: AR-C51-shar, AR-C51-qadam, AR-C51-isba, AR-C51-katif, AR-C51-raqaba

\chapter{Hair, Foot, Finger, Shoulder, Neck}
\label{ch:ar-hair-foot-finger-shoulder-neck}
\hlchaptermodality{51}

\begin{chapteropening}
\textbf{By the end of this chapter:} \emph{I can say hair, a foot, a finger, a shoulder and the neck.}

Complete the last lesson of chapter 51: I can say hair, a foot, a finger, a shoulder and the neck.
\end{chapteropening}

\section[\texorpdfstring{\ar{شعر} (shaʿr)}{شعر (shaʿr)}]{\ar{شعر} (shaʿr) — hair}

\label{lesson:AR-C51-shar}



\begin{quote}
Before the new one: say the Arabic for a mouth, then the Arabic for a tooth.
\end{quote}

\subsection*{You'll want to know: \ar{شعر}}
\textbf{\ar{شعر}} — \emph{shaʿr} — \textquotedblleft{}hair\textquotedblright{}.

\begin{cousinweb}
\textbf{\ar{شعر}} (\emph{shaʿr}) is built on the native Arabic root \emph{sh-ʿ-r}, \textquotedblleft{}hair; to perceive\textquotedblright{} (Lane). The same root gives \emph{shiʿr}, \textquotedblleft{}poetry\textquotedblright{}; Hebrew \emph{seʿar}, \textquotedblleft{}hair\textquotedblright{}, is usually counted a cousin.
\end{cousinweb}

\begin{grammarlens}[title={What you've built}]
One more word for this chapter.
\end{grammarlens}

\subsection*{Your turn}
\begin{itemize}
\raggedright
  \item \emph{Say it:} \emph{shaʿr}
  \item \emph{Say it:} \emph{shaʿr}, once more
  \item \emph{Say it:} say \emph{shaʿr}
  \item \emph{Recall:} say the Arabic for a mouth, then the Arabic for a tooth, then say \emph{shaʿr} again
\end{itemize}

\begin{tcolorbox}[breakable,colback=teal!4,colframe=teal!35!black,title={Before you move on}]
What does \ar{شعر} mean? (\textquotedblleft{}Hair\textquotedblright{}.) Say it once more.
\end{tcolorbox}
\section[\texorpdfstring{\ar{قدم} (qadam)}{قدم (qadam)}]{\ar{قدم} (qadam) — a foot}

\label{lesson:AR-C51-qadam}



\begin{quote}
Before the new one: say the Arabic for a tooth, then the Arabic for hair.
\end{quote}

\subsection*{You'll want to know: \ar{قدم}}
\textbf{\ar{قدم}} — \emph{qadam} — \textquotedblleft{}a foot\textquotedblright{}.

\begin{cousinweb}
\textbf{\ar{قدم}} (\emph{qadam}) is built on the native Arabic root \emph{q-d-m}, \textquotedblleft{}to go before, to precede\textquotedblright{} (Lane). \emph{Qadam} is the foot, the part that goes first; the same root gives \emph{qadīm}, \textquotedblleft{}old\textquotedblright{}.
\end{cousinweb}

\begin{grammarlens}[title={What you've built}]
One more word for this chapter.
\end{grammarlens}

\subsection*{Your turn}
\begin{itemize}
\raggedright
  \item \emph{Say it:} \emph{qadam}
  \item \emph{Say it:} \emph{qadam}, once more
  \item \emph{Say it:} say \emph{qadam}
  \item \emph{Recall:} say the Arabic for a tooth, then the Arabic for hair, then say \emph{qadam} again
\end{itemize}

\begin{tcolorbox}[breakable,colback=teal!4,colframe=teal!35!black,title={Before you move on}]
What does \ar{قدم} mean? (\textquotedblleft{}A foot\textquotedblright{}.) Say it once more.
\end{tcolorbox}
\section[\texorpdfstring{\ar{إصبع} (iṣbaʿ)}{إصبع (iṣbaʿ)}]{\ar{إصبع} (iṣbaʿ) — a finger}

\label{lesson:AR-C51-isba}



\begin{quote}
Before the new one: say the Arabic for hair, then the Arabic for a foot.
\end{quote}

\subsection*{You'll want to know: \ar{إصبع}}
\textbf{\ar{إصبع}} — \emph{iṣbaʿ} — \textquotedblleft{}a finger\textquotedblright{}.

\begin{cousinweb}
\textbf{\ar{إصبع}} (\emph{iṣbaʿ}) is built on the native Arabic root \emph{ṣ-b-ʿ}, \textquotedblleft{}a finger\textquotedblright{} (Lane). Hebrew \emph{etsbaʿ}, \textquotedblleft{}finger\textquotedblright{}, is its cousin.
\end{cousinweb}

\begin{grammarlens}[title={What you've built}]
One more word for this chapter.
\end{grammarlens}

\subsection*{Your turn}
\begin{itemize}
\raggedright
  \item \emph{Say it:} \emph{iṣbaʿ}
  \item \emph{Say it:} \emph{iṣbaʿ}, once more
  \item \emph{Say it:} say \emph{iṣbaʿ}
  \item \emph{Recall:} say the Arabic for hair, then the Arabic for a foot, then say \emph{iṣbaʿ} again
\end{itemize}

\begin{tcolorbox}[breakable,colback=teal!4,colframe=teal!35!black,title={Before you move on}]
What does \ar{إصبع} mean? (\textquotedblleft{}A finger\textquotedblright{}.) Say it once more.
\end{tcolorbox}
\section[\texorpdfstring{\ar{كتف} (katif)}{كتف (katif)}]{\ar{كتف} (katif) — a shoulder}

\label{lesson:AR-C51-katif}



\begin{quote}
Before the new one: say the Arabic for a foot, then the Arabic for a finger.
\end{quote}

\subsection*{You'll want to know: \ar{كتف}}
\textbf{\ar{كتف}} — \emph{katif} — \textquotedblleft{}a shoulder\textquotedblright{}.

\begin{cousinweb}
\textbf{\ar{كتف}} (\emph{katif}) is built on the native Arabic root \emph{k-t-f}, \textquotedblleft{}the shoulder, the shoulder-blade\textquotedblright{} (Lane). Hebrew \emph{katef}, \textquotedblleft{}shoulder\textquotedblright{}, is its cousin.
\end{cousinweb}

\begin{grammarlens}[title={What you've built}]
One more word for this chapter.
\end{grammarlens}

\subsection*{Your turn}
\begin{itemize}
\raggedright
  \item \emph{Say it:} \emph{katif}
  \item \emph{Say it:} \emph{katif}, once more
  \item \emph{Say it:} say \emph{katif}
  \item \emph{Recall:} say the Arabic for a foot, then the Arabic for a finger, then say \emph{katif} again
\end{itemize}

\begin{tcolorbox}[breakable,colback=teal!4,colframe=teal!35!black,title={Before you move on}]
What does \ar{كتف} mean? (\textquotedblleft{}A shoulder\textquotedblright{}.) Say it once more.
\end{tcolorbox}
\section[\texorpdfstring{\ar{رقبة} (raqaba)}{رقبة (raqaba)}]{\ar{رقبة} (raqaba) — the neck}

\label{lesson:AR-C51-raqaba}



\begin{quote}
Before the new one: say the Arabic for a finger, then the Arabic for a shoulder.
\end{quote}

\subsection*{You'll want to know: \ar{رقبة}}
\textbf{\ar{رقبة}} — \emph{raqaba} — \textquotedblleft{}the neck\textquotedblright{}.

\begin{cousinweb}
\textbf{\ar{رقبة}} (\emph{raqaba}) is built on the native Arabic root \emph{r-q-b}, \textquotedblleft{}the neck; to watch over\textquotedblright{} (Lane). The verb \emph{raqaba}, \textquotedblleft{}to watch\textquotedblright{}, sits under the same root in Lane.
\end{cousinweb}

\begin{grammarlens}[title={What you've built}]
One more word for this chapter.
\end{grammarlens}

\subsection*{Your turn}
\begin{itemize}
\raggedright
  \item \emph{Say it:} \emph{raqaba}
  \item \emph{Say it:} \emph{raqaba}, once more
  \item \emph{Say it:} say \emph{raqaba}
  \item \emph{Recall:} say the Arabic for a finger, then the Arabic for a shoulder, then say \emph{raqaba} again
\end{itemize}

\begin{tcolorbox}[breakable,colback=teal!4,colframe=teal!35!black,title={Before you move on}]
What does \ar{رقبة} mean? (\textquotedblleft{}The neck\textquotedblright{}.) Say it once more.
\end{tcolorbox}
